\section{图神经网络的泛函算子解释}

\label{sec:12-2}

图神经网络最常见的工程描述是“聚合邻居再更新表示”，但若只停留在这个层面，很多关键结构都会被遮蔽：为什么图卷积总要与拉普拉斯或归一化邻接联系在一起，为什么平滑、扩散和低通滤波会反复出现，以及哪些图学习步骤真的能够被第 \ref{chap:09} 章和第 \ref{chap:10} 章的算子语言准确接口化。本节选择从图拉普拉斯与 Dirichlet 能量出发，抓住那些可被谱论、prox 和分裂方法稳定描述的部分。

\subsection{图信号、拉普拉斯与能量}

\begin{definition}[图拉普拉斯算子]\label{def:graph-laplacian-operator}
设 $\mathscr G=(V,E,w)$ 为有限加权无向图，其中 $V=\{1,\dots,n\}$，边权满足 $w_{ij}=w_{ji}\ge 0$。记图信号空间
\[
H_V:=\mathbb R^n
\]
并赋予标准内积。定义图拉普拉斯算子 $L_{\mathscr G}\colon H_V\to H_V$ 为
\[
(L_{\mathscr G}u)_i:=\sum_{j=1}^n w_{ij}(u_i-u_j),
\qquad i=1,\dots,n.
\]
\end{definition}

\begin{definition}[图 Dirichlet 能量]\label{def:graph-dirichlet-energy}
在定义~\ref{def:graph-laplacian-operator} 的条件下，称泛函
\[
\mathcal E_{\mathscr G}(u)
:=
\frac12\sum_{i,j=1}^n w_{ij}(u_i-u_j)^2
\]
为图 $\mathscr G$ 上信号 $u$ 的\emph{Dirichlet 能量}。
\end{definition}

\begin{proposition}[图拉普拉斯的正性]\label{prop:graph-laplacian-positive}
算子 $L_{\mathscr G}$ 是 $H_V$ 上的自伴正半定算子，并满足
\[
\langle L_{\mathscr G}u,u\rangle=\mathcal E_{\mathscr G}(u),
\qquad \forall u\in H_V.
\]
\end{proposition}

\begin{proof}
由权重对称性，
\[
\langle L_{\mathscr G}u,v\rangle
=
\sum_{i=1}^n \sum_{j=1}^n w_{ij}(u_i-u_j)v_i
=
\frac12\sum_{i,j=1}^n w_{ij}(u_i-u_j)(v_i-v_j),
\]
交换 $u,v$ 可知它自伴。令 $v=u$，得到
\[
\langle L_{\mathscr G}u,u\rangle
=
\frac12\sum_{i,j=1}^n w_{ij}(u_i-u_j)^2
=
\mathcal E_{\mathscr G}(u)\ge 0.
\]
故 $L_{\mathscr G}$ 正半定。
\end{proof}

\begin{remark}
由于 $H_V$ 是有限维 Hilbert 空间，$L_{\mathscr G}$ 同时还是紧自伴算子。于是第 \ref{sec:3-4} 节定理~\ref{thm:compact-selfadjoint-spectral} 可直接应用于它：图傅里叶基、特征值平滑强度和谱滤波，本质上都是紧自伴谱理论在图上的有限维坍缩。
\end{remark}

\subsection{消息传递与图滤波}

\begin{proposition}[消息传递是图滤波的一种线性原型]\label{prop:message-passing-graph-filter}
设 $p$ 为一元多项式，则算子
\[
T_p:=p(L_{\mathscr G})
\]
与 $L_{\mathscr G}$ 共享同一组特征向量，因此是一个图谱滤波器。特别地，线性消息传递层
\[
u^{+}=(I-\tau L_{\mathscr G})u,
\qquad 0<\tau<\|L_{\mathscr G}\|^{-1},
\]
对应于一阶低通图滤波；若再附加逐点非线性与通道权矩阵，就得到常见 GNN 层的线性骨架。
\end{proposition}

\begin{proof}
由定理~\ref{thm:compact-selfadjoint-spectral}，存在正交基 $\{e_k\}_{k=1}^n$ 与实特征值 $\lambda_k\ge 0$，使
\[
L_{\mathscr G}e_k=\lambda_k e_k.
\]
于是
\[
T_p e_k=p(\lambda_k)e_k,
\]
故 $T_p$ 与 $L_{\mathscr G}$ 同时对角化。特别地，当 $p(\lambda)=1-\tau\lambda$ 时，
\[
T_p=I-\tau L_{\mathscr G},
\]
它衰减高频特征模态，正是一阶低通滤波。
\end{proof}

\begin{theorem}[图扩散与学习接口]\label{thm:graph-diffusion-learning-interface}
设 $\tau>0$，$y\in H_V$。则问题
\[
\min_{u\in H_V}
\left\{
\frac12\|u-y\|^2+\tau \mathcal E_{\mathscr G}(u)
\right\}
\]
有唯一极小元 $u_\tau$，并满足
\[
u_\tau=(I+\tau L_{\mathscr G})^{-1}y=\operatorname{prox}_{\tau\mathcal E_{\mathscr G}}(y).
\]
此外，映射
\[
y\longmapsto (I-\tau L_{\mathscr G})y
\]
是上述 resolvent 的一阶显式近似，因此线性图扩散、图平滑和一类 GNN 消息传递层都可被看成第 \ref{sec:10-1} 节定义~\ref{def:forward-backward-iteration} 在图信号空间上的具体接口。
\end{theorem}

\begin{proof}
由命题~\ref{prop:graph-laplacian-positive}，
\[
\mathcal E_{\mathscr G}(u)=\langle L_{\mathscr G}u,u\rangle
\]
是凸二次型。故目标是严格凸连续函数，唯一极小元存在。对极小元 $u_\tau$ 写一阶最优性条件，
\[
0=u_\tau-y+\tau L_{\mathscr G}u_\tau.
\]
于是
\[
(I+\tau L_{\mathscr G})u_\tau=y,
\qquad
u_\tau=(I+\tau L_{\mathscr G})^{-1}y.
\]
另一方面，由定义~\ref{def:proximal-mapping}，$\operatorname{prox}_{\tau\mathcal E_{\mathscr G}}(y)$ 正是上式的唯一极小元，故
\[
u_\tau=\operatorname{prox}_{\tau\mathcal E_{\mathscr G}}(y).
\]
最后，由 resolvent 展开
\[
(I+\tau L_{\mathscr G})^{-1}y
=
y-\tau L_{\mathscr G}y+O(\tau^2),
\]
可见显式一步 $y\mapsto y-\tau L_{\mathscr G}y$ 是对应的前向近似。
\end{proof}

\begin{remark}
定理~\ref{thm:graph-diffusion-learning-interface} 只 formalize 那些可嵌入凸能量或算子分裂语言的图学习部分。一般深层 GNN 训练包含参数共享、非线性组合与非凸优化；本节并不把这些内容伪装成凸问题，而是只抽取其中最稳定的图滤波、平滑与扩散骨架。
\end{remark}

\subsection{典型例子}

\begin{example}[图平滑最小化]
若 $y$ 是带噪观测信号，则定理~\ref{thm:graph-diffusion-learning-interface} 给出的极小元平衡了数据保真和图平滑。这正是半监督学习、标签传播与图正则回归中的基础模型。
\end{example}

\begin{example}[邻接矩阵与归一化拉普拉斯的有限维坍缩]
在有限图上，常见 GCN 层写作
\[
H^{+}=\sigma(\widetilde D^{-1/2}\widetilde A\widetilde D^{-1/2}HW).
\]
把
\[
\widetilde L:=I-\widetilde D^{-1/2}\widetilde A\widetilde D^{-1/2}
\]
代入，就得到
\[
\widetilde D^{-1/2}\widetilde A\widetilde D^{-1/2}=I-\widetilde L,
\]
因此消息传递矩阵本质上就是一阶图滤波器。
\end{example}

\begin{example}[知识图谱与 item-item 推荐中的关系传播]
在知识图谱推荐或 item-item 推荐中，边权编码相似度、共现或关系强度。若只保留图传播的线性骨架，则每一步更新都可解释为对邻域一致性和局部平滑的近似求解。把这个例子放在这里，是为了把 GNN 直接接到推荐系统的图结构接口上。
\end{example}

\subsection*{有限维对照}

Boyd 不会直接讲 GNN，但他关于二次型、谱分解和分布式平滑的有限维母语，在这里仍然有效。区别只在于：工程上常把“邻接矩阵乘一次”当成一个架构模块，而本节说明它其实是图拉普拉斯谱滤波、Dirichlet 能量最小化以及 prox/前向步的一种坍缩实现。理解这一点以后，第 \ref{sec:12-4} 节的分布式图推荐模型就不再只是经验配方，而能重新接回 ADMM 与原--对偶分裂。

\subsection*{习题（待补）}

\begin{exercise}[图拉普拉斯正性占位]
补一题：证明命题~\ref{prop:graph-laplacian-positive}，并说明常数信号为什么总落在零特征值对应空间中。
\end{exercise}

\begin{exercise}[图扩散与 prox 占位]
补一题：从定理~\ref{thm:graph-diffusion-learning-interface} 出发，说明图扩散为何可以被写成定义~\ref{def:proximal-mapping} 的邻近算子。
\end{exercise}
